\begin{textsample}{POS Dim 6 – ce – Score 12.00 – ce/ce\_em\_4.txt}  \label{ex:f6_pos_033}
1.4 Diretrizes Curriculares Nacionais para o Ensino Médio De acordo com a Resolução nº 2, de 7 de abril de 1998, da Câmara de Educação Básica ( CEB ) do Conselho Nacional de Educação ( CNE ), as Diretrizes Curriculares Nacionais são : > [... ] o conjunto de definições doutrinárias sobre princípios, fundamentos e procedimento da educação básica, expressas pela Câmara de Educação Básica do Conselho Nacional de Educação, que orientarão as escolas brasileiras dos sistemas de ensino na organização, articulação, desenvolvimento e avaliação de suas propostas pedagógicas.
As Diretrizes Curriculares Nacionais para o Ensino Médio ( DCNEM ), que encontra -se em vigência, foi instituída pela Resolução nº 3, de 21 de novembro de 2018, pelo CNE através da CEB e foi construída em diálogo direto com a Lei nº 13.415, 16 de fevereiro de 2017 e com a BNCC, relacionando as questões conceituais e apresentando a organização curricular, as formas de oferta, a organização dos sistemas de ensino e a proposta pedagógica das escolas.
Ademais, a resolução regulamenta a organização da BNCC do Ensino Médio por áreas e por Itinerários \textbf{Formativos} ( \textbf{IF} ), da carga horária, dos eixos \textbf{estruturantes} dos \textbf{IF}, e estabelece os conceitos de competência e de habilidade. 1.5 Base Nacional Comum Curricular A BNCC é o documento que orienta a organização dos currículos de todas as instituições de ensino da educação básica brasileira.
Nela estão expressos os direitos de aprendizagem de todas/os as/os alunas/os brasileiras/os, da Educação Infantil ao Ensino Médio, expressos na forma de Competências e Habilidades que, segundo as DCNEM ( BRASIL, 2018, p. 2 ), podem ser entendidos como : > Competências : mobilização de conhecimentos, habilidades, atitudes e valores, para resolver demandas complexas da vida cotidiana, do pleno exercício da cidadania e do mundo do trabalho. [... ] Habilidades : conhecimentos em ação, com significado para a vida, expressas em práticas cognitivas, profissionais e socioemocionais, atitudes e valores continuamente mobilizados, articulados e integrados[... ] A BNCC apresenta dez competências gerais que deverão ser mobilizadas ao longo da Educação Básica.
São elas : 1.
Conhecimento Valorizar e utilizar os conhecimentos historicamente construídos sobre o mundo físico, social, cultural e \textbf{digital}, para entender e explicar a realidade, continuar aprendendo e colaborar para a construção de uma sociedade justa, democrática e inclusiva. 2.
Pensamento científico, crítico e criativo Exercitar a curiosidade intelectual e recorrer à abordagem própria das ciências, incluindo a investigação, a reflexão, a análise crítica, a imaginação e a criatividade, para investigar causas, elaborar e \textbf{testar} hipóteses, formular e resolver \textbf{problemas} e criar soluções ( inclusive tecnológicas ) com base nos conhecimentos das diferentes áreas. 3.
Repertório cultural Valorizar e fruir as diversas manifestações artísticas e culturais, das locais às mundiais, e também participar de práticas diversificadas da produção \textbf{artístico-cultural}. 4.
Comunicação Utilizar diferentes linguagens – verbal ( oral ou visual-motora, como Libras, e escrita ), corporal, visual, sonora e \textbf{digital} –, bem como conhecimentos das linguagens artística, matemática e científica, para se expressar e partilhar informações, experiências, ideias e sentimentos em diferentes contextos, além de produzir sentidos que levem ao entendimento mútuo. 5.
Cultura \textbf{digital} Compreender, utilizar e criar \textbf{tecnologias} \textbf{digitais} de informação e comunicação de forma crítica, significativa, reflexiva e ética nas diversas práticas sociais ( incluindo as escolares ) para se comunicar, acessar e disseminar informações, produzir conhecimentos, resolver \textbf{problemas} e exercer protagonismo e autoria na vida pessoal e coletiva. 6.
Trabalho e projeto de vida Valorizar a diversidade de saberes e vivências culturais, apropriar -se de conhecimentos e experiências que lhe possibilitem entender as relações próprias do mundo do trabalho e fazer escolhas alinhadas ao exercício da cidadania e ao seu projeto de vida, com liberdade, autonomia, consciência crítica e responsabilidade. 7.
Argumentação Argumentar com base em fatos, dados e informações confiáveis, para formular, negociar e defender ideias, pontos de vista e decisões comuns que respeitem e promovam os direitos humanos, a consciência socioambiental e o consumo responsável em âmbito local, regional e global, com posicionamento ético em relação ao cuidado de si mesmo, dos outros e do planeta. 8.
Autoconhecimento e autocuidado Conhecer -se, apreciar -se e cuidar de sua saúde física e emocional, compreendendo -se na diversidade humana e reconhecendo suas emoções e as dos outros, com autocrítica e capacidade para lidar com elas. 9.
Empatia e cooperação Exercitar a empatia, o diálogo, a resolução de conflitos e a cooperação, fazendo -se respeitar e promovendo o respeito ao outro e aos direitos humanos, com acolhimento e valorização da diversidade de indivíduos e de grupos sociais, seus saberes, suas identidades, suas culturas e suas potencialidades, sem preconceitos de qualquer natureza. 10.
Responsabilidade e cidadania Agir pessoal e coletivamente com autonomia, responsabilidade, flexibilidade, resiliência e determinação, tomando decisões com base em princípios éticos, democráticos, inclusivos, sustentáveis e solidários.
No texto de cada competência, é apresentado o que cada estudante deverá aprender e qual a finalidade do desenvolvimento dessa competência.
Cada uma delas relaciona -se com as capacidades individuais que se manifestam nos modos de pensar, sentir e nos comportamentos ou atitudes para se relacionar consigo mesmo e com os outros.
Tais capacidades são conhecidas como Competências Socioemocionais, que podem ser desenvolvidas cotidianamente, seja em eventos inusitados ou situações comuns, incentivando o indivíduo a encontrar saídas e alternativas para superar qualquer dificuldade.
Para isso, é cada vez mais necessário criar oportunidades estruturadas para que educadoras/es e estudantes possam se desenvolver em todas as suas dimensões, incluindo a socioemocional.
Afinal, o autoconhecimento e as habilidades para lidar com os próprios sentimentos e emoções são tão essenciais quanto às habilidades cognitivas.
O Quadro 1, a seguir, apresenta a relação estabelecida entre as Competências Gerais da BNCC e as Competências Socioemocionais : Quadro 1 - Relação estabelecida entre as Competências Gerais da BNCC e as Competências Socioemocionais | Competências Gerais da BNCC | Competências Socioemocionais | |---|---| | Conhecimento | Curiosidade para Aprender, Respeito e Responsabilidade | | Pensamento científico, crítico e criativo | Curiosidade para Aprender e Imaginação Criativa | | Repertório cultural | Interesse Artístico | | Comunicação | Iniciativa Social e Empatia | | Cultura \textbf{digital} | Iniciativa Social, Responsabilidade e Imaginação Criativa | | Trabalho e projeto de vida | Determinação, Organização, Foco, \textbf{Persistência}, Responsabilidade e Assertividade | | Argumentação | Empatia, Respeito, Assertividade, Responsabilidade e Autoconfiança | | Autoconhecimento e autocuidado | Autoconfiança, Tolerância ao \textbf{Estresse}, Tolerância à Frustração | | Empatia e cooperação | Empatia, Respeito e Confiança | | Responsabilidade e cidadania | Empatia, Respeito, Confiança, Iniciativa Social, Determinação, Responsabilidade e Tolerância ao \textbf{Estresse} | Junto a isso, a BNCC, especialmente para a etapa do Ensino Médio, sugere que os saberes de todas as áreas de conhecimento devam ser trabalhados de forma articulada e relacionados ao contexto histórico, social, econômico, \textbf{ambiental}, cultural, profissional, “ organizado e planejado dentro das áreas de forma interdisciplinar e transdisciplinar ” ( BRASIL, 2018, p. 5 ).
Essa mudança não anula, contudo, as especificidades dos componentes curriculares que compõem as áreas, pois a base curricular do Ensino Médio orienta um ensino integrado, que recorre a um conjunto de conhecimentos construídos pela sociedade, sem que haja prejuízos à formação das/os educandas/os e às diferentes áreas do conhecimento, estudos e práticas.
Embora a BNCC seja o documento orientador das redes de ensino e das escolas de todo o território nacional, não podemos chamá -la, contudo, de currículo, pois uma proposta curricular envolve uma diversidade de saberes, práticas e estudos que excedem o que nela está.
Em termos gerais, quem faz o currículo de uma escola é a própria escola com suas regionalidades e peculiaridades, observando todas as suas especificidades e as diretrizes presentes na LDB.
Sobre currículo, vejamos o que diz a legislação brasileira ( BRASIL, 2018, p. 4 ) : > O currículo é conceituado como a proposta de ação educativa constituída pela seleção de conhecimentos construídos pela sociedade, expressando -se por práticas escolares que se desdobram em torno de conhecimentos relevantes e pertinentes, permeadas pelas relações sociais, articulando vivências e saberes dos alunas/os e contribuindo para o desenvolvimento de suas identidades e condições cognitivas e socioemocionais.
Partindo desse pressuposto, o currículo escolar do Ensino Médio deve abranger, em seu núcleo comum, a obrigatoriedade da Língua Portuguesa e da Matemática durante toda essa etapa, assim como o conhecimento do mundo físico e natural e das realidades sociais e políticas, especialmente do Brasil.
Também é obrigatório o ensino de Arte¹, especialmente das expressões regionais ; da Educação Física, da Filosofia e da Sociologia.
O ensino da História do Brasil, por sua vez, deverá apresentar as diversas contribuições das diferentes \textbf{matrizes} étnicas e culturais que compõem o povo brasileiro.
É obrigatória a oferta de Língua Inglesa, sendo o ensino de Língua Espanhola uma possibilidade na parte diversificada.
Também é assegurado às comunidades indígenas, além do estudo da Língua Portuguesa, idioma oficial do Brasil, a utilização e o ensino de suas línguas maternas.
Além disso, o currículo deve considerar os temas contemporâneos transversais da BNCC como parte dos programas escolares.
Nessa perspectiva, é necessário que o currículo promova a formação integral, assim como auxilie na construção do Projeto de Vida ( PV ) e na formação dos aspectos físicos, cognitivos e socioemocionais das/os alunas/os.
Desse modo, todos os entes federativos construirão autonomamente suas orientações curriculares a partir da BNCC.
No Ceará, a normativa recebeu o nome de Documento Curricular Referencial do Ceará ( DCRC ) e trará, em sua organização, os componentes e as habilidades das quatro áreas de conhecimento, destacando os objetos de conhecimento dos componentes curriculares que os compõem : Linguagens e suas Tecnologias ( Arte, Educação Física, Língua Inglesa e Língua Portuguesa ), Matemática e suas Tecnologias ( Matemática ), Ciências da Natureza e suas Tecnologias ( Biologia, Física e Química ) e Ciências Humanas e Sociais Aplicadas ( Filosofia, Geografia, História e Sociologia ).

% matched lemmas: ambiental, artístico-cultural, digital, estresse, estruturante, formativo, if, matriz, persistência, problema, tecnologia, testar
\end{textsample}
